
%BeginMC
\question

\begin{choices}
\correctchoice  %EndChoice
\choice  %EndChoice
\choice  %EndChoice
\choice  %EndChoice
\choice  %EndChoice
\end{choices}
%EndMC


%=============

%BeginMC
\question
A marine iguana is basking on a sunny rock after feeding in the cold Pacific Ocean water. It's warming its body through 

\begin{choices}
\correctchoice conduction and radiation. %EndChoice
\choice convection and evaporation. %EndChoice
\choice conduction and convection %EndChoice
\choice convection and radiation.  %EndChoice
\choice conduction and evaporation. %EndChoice
\end{choices}
%EndMC


%BeginMC
\question
Stomach cells are moderately well adapted to the acidity and digesting activities in the stomach by having

\begin{choices}
\correctchoice a thick, mucous secretion and active replacement of epithelial cells. %EndChoice
\choice a sufficient colony of \textit{H. pylori} to breakdown the cellulose from vegetables. %EndChoice
\choice a high level of secretion of pepsin. %EndChoice
\choice a high level of bile. %EndChoice
\choice secretions enter the stomach from the pancreas. %EndChoice
\end{choices}
%EndMC


%BeginMC
\question
Upon activation by stomach acidity, the secretions of the gastric glands

\begin{choices}
\correctchoice initiate the digestion of protein in the stomach. %EndChoice
\choice initiate the mechanical digestion of lipids in the stomach. %EndChoice
\choice initiate the chemical digestion of lipids in the stomach.  %EndChoice
\choice include nitrogenase and amylase. %EndChoice
\choice delay digestion until the food arrives in the small intestine. %EndChoice
\end{choices}
%EndMC



%BeginMC
\question
The osmoregulatory/excretory system of a freshwater flatworm is based on the operation of

\begin{choices}
\correctchoice protonephridia. %EndChoice
\choice metanephridia. %EndChoice
\choice Malpighian tubules. %EndChoice
\choice nephrons. %EndChoice
\choice ananephredia. %EndChoice
\end{choices}
%EndMC


%BeginMC
\question
The transfer of fluid from the glomerulus to Bowman's capsule

\begin{choices}
\correctchoice  is mainly a consequence of blood pressure in the capillaries of the glomerulus. %EndChoice
\choice results from active transport. %EndChoice
\choice transfers large molecules  as easily as small ones. %EndChoice
\choice is very selective as to which protein-sized molecules are transferred. %EndChoice
\choice usually includes the transfer of red blood cells into Bowman's capsule. %EndChoice
\end{choices}
%EndMC


%BeginMC
\question
Which of the following is an accurate statement about thermoregulation?

\begin{choices}
\correctchoice Endotherms and ectotherms differ in their primary source of heat for thermoregulation. %EndChoice
\choice Endotherms are regulators and ectotherms are conformers. %EndChoice
\choice Endotherms maintain a constant body temperature and ectotherms do not. %EndChoice
\choice Endotherms are warm-blooded and ectotherms are cold-blooded.  %EndChoice
\choice Endothermy has a lower energy cost than ectothermy. %EndChoice
\end{choices}
%EndMC



%BeginMC
\question
The advantage of excreting nitrogenous wastes as urea rather than as ammonia is that

\begin{choices}
\correctchoice urea is less toxic than ammonia. %EndChoice
\choice urea requires more water for excretion than ammonia. %EndChoice
\choice urea can be exchanged for \ce{Na+}. %EndChoice
\choice urea does not affect the osmoregulation. %EndChoice
\choice less nitrogen is removed from the body. %EndChoice
\end{choices}
%EndMC


%BeginMC
\question
Birds secrete uric acid as their nitrogenous waste because uric acid

\begin{choices}
\correctchoice requires little water for nitrogenous waste disposal, thus reducing body mass. %EndChoice
\choice is readily soluble in water. %EndChoice
\choice is metabolically less expensive to synthesize than other excretory products. %EndChoice
\choice excretion allows birds to live in desert environments. %EndChoice
\choice is less toxic than urea. %EndChoice
\end{choices}
%EndMC

%BeginMC
\question
Stomata will open when  

\begin{choices}
\correctchoice the guard cells have a higher solute concentration than outside the cell. %EndChoice
\choice the guard cells have a lower solute concentration than outside the cell. %EndChoice
\choice the solute potential inside the guard cell is positive. %EndChoice
\choice the pressure potential inside the guard cell is negative. %EndChoice
\choice the water potential inside the guard cell is higher than outside the cell. %EndChoice
\end{choices}
%EndMC


%BeginMC
\question
Which of the following types of molecules are the major structural components of the cell membrane?

\begin{choices}
\choice phospholipids and cellulose. %EndChoice
\choice nucleic acids and proteins. %EndChoice
\correctchoice phospholipids and proteins. %EndChoice
\choice proteins and cellulose. %EndChoice
\choice proteins and collagen. %EndChoice
\end{choices}
%EndMC

%BeginMC
\question
Which of the following contains digestive enzymes?

\begin{choices}
\correctchoice lysosome %EndChoice
\choice food vacuole. %EndChoice
\choice central vacuole. %EndChoice
\choice guard cells. %EndChoice
\choice contractile vacuole. %EndChoice
\end{choices}
%EndMC




%BeginMC
\question
Which of the following statements is correct about diffusion in or out of a cell?

\begin{choices}
\choice It requires aquaporins in the cell wall. %EndChoice
\choice It requires an expenditure of energy by the cell. %EndChoice
\correctchoice It is a passive process in which molecules move from a region of higher concentration to a region of lower concentration. %EndChoice
\choice It is an active process in which molecules move from a region of lower concentration to one of higher concentration. %EndChoice
\choice It requires channel proteins in the cell membrane. %EndChoice
\end{choices}
%EndMC


%BeginMC
\question
Which of the following membrane activities requires energy from ATP?

\begin{choices}
\choice diffusion of chloride ions across the membrane through a chloride channel. %EndChoice
\choice movement of water into a cell. %EndChoice
\correctchoice movement of \ce{Na+} ions from a lower concentration in a mammalian cell to a higher concentration in the extracellular fluid. %EndChoice
\choice movement of glucose molecules into a bacterial cell from a solution containing a higher concentration of glucose than inside the cell. %EndChoice
\choice movement of water out of a paramecium. %EndChoice
\end{choices}
%EndMC


%BeginMC
\question
A species of rabbit that follows Bergmann's rule 

\begin{choices}
\correctchoice  would be larger in Canada than in Texas. %EndChoice
\choice would be smaller in Minnesota than in Florida. %EndChoice
\choice would be whiter in winter and browner in summer. %EndChoice
\choice would have larger ears in Maine and smaller ears in Georgia. %EndChoice
\end{choices}
%EndMC

%BeginMC
\question
Small mammals like mice use more energy per unit mass than large mammals mostly because most of the energy is used for

\begin{choices}
\correctchoice thermoregulation. %EndChoice
\choice the basal metabolic rate. %EndChoice
\choice activities like hunting. %EndChoice
\choice reproduction. %EndChoice
\choice growing quickly. %EndChoice
\end{choices}
%EndMC


%BeginMC
\question
Cell walls are found in at least some representatives of  
\begin{choices}
\choice algae %EndChoice
\choice plants %EndChoice
\choice fungi %EndChoice
\choice protists %EndChoice
\choice bacteria %EndChoice
\correctchoice all of the above.  %EndChoice %dontshuffle
\end{choices}
%EndMC



%BeginMC
\question
A fungus is a type of 
\begin{choices}
\correctchoice chemoheterotroph %EndChoice
\choice chemoautotroph %EndChoice
\choice  photoheterotroph %EndChoice
\choice  photoautotroph %EndChoice
\choice  autoheterotroph %EndChoice
\end{choices}
%EndMC

%BeginMC
\question
Organisms that obtain energy by oxidizing iron compounds and obtain carbon from \ce{CO2} are called

\begin{choices}
\correctchoice chemoautotrophs. %EndChoice
\choice chemoheterotrophs. %EndChoice
\choice photoheterotrophs. %EndChoice
\choice photoautotrophs. %EndChoice
\choice autoheterotrophs. %EndChoice
\choice photochemotrophs. %EndChoice
\end{choices}
%EndMC

%BeginMC
\question
Ingested dietary substances must cross cell membranes to be used by the body, a process known as

\begin{choices}
\choice ingestion. %EndChoice
\choice digestion. %EndChoice
\choice hydrolysis. %EndChoice
\correctchoice absorption. %EndChoice
\choice elimination. %EndChoice
\end{choices}
%EndMC

%BeginMC
\question
Because the foods eaten by animals are often composed largely of macromolecules, this requires the animals to have mechanisms for \rule{1in}{0.4pt} not found in other groups.

\begin{choices}
\choice elimination of nitrogenous waste. %EndChoice
\choice absorption of nutrients. %EndChoice
\correctchoice chemical digestion. %EndChoice
\choice reabsorption of solutes from the loop of Henle. %EndChoice
\choice filtration in Bowman's capsule. %EndChoice
\end{choices}
%EndMC


%BeginMC
\question
After ingestion by humans, the first category of macromolecules to be chemically digested by enzymes in the mouth is

\begin{choices}
\choice proteins. %EndChoice
\correctchoice carbohydrates. %EndChoice
\choice cholesterol and other lipids. %EndChoice
\choice nucleic acids. %EndChoice
\choice minerals. %EndChoice
\end{choices}
%EndMC

%BeginMC
\question
Pepsin is a digestive enzyme that

\begin{choices}
\choice secreted by the pancreas to digest proteins in the duodenum. %EndChoice
\choice present in bile to digest lipids in the small intestine. %EndChoice
\choice digests polysaccharides to monosaccharides in the small intestine. %EndChoice
\correctchoice begins the digestion of proteins in the stomach. %EndChoice
\choice digests nucleic acids in the stomach. %EndChoice
\end{choices}
%EndMC

%BeginMC
\question
When the digestion and absorption of carbohydrates results in more energy-rich molecules than are immediately required by an animal, the excess is

\begin{choices}
\choice eliminated in the feces. %EndChoice
\choice stored as starch. %EndChoice
\correctchoice stored as glycogen or fat. %EndChoice
\choice stored in the crop. %EndChoice
\choice stored as cellulose or chitin. %EndChoice
\end{choices}
%EndMC


%BeginMC
\question
Which of the following organs is incorrectly paired with its function?

\begin{choices}
\choice Stomach-protein digestion. %EndChoice
\choice Mouth-starch digestion. %EndChoice
\correctchoice Large intestine-bile production. %EndChoice
\choice Small intestine-nutrient absorption. %EndChoice
\choice Pancreas-enzyme production. %EndChoice
\end{choices}
%EndMC

%BeginMC
\question
Which of the following is \emph{not} a major activity of the stomach?

\begin{choices}
\choice Mechanical digestion. %EndChoice
\choice \ce{HCl} secretion. %EndChoice
\choice Mucus secretion. %EndChoice
\correctchoice Nutrient absorption. %EndChoice
\choice Enzyme secretion. %EndChoice
\end{choices}
%EndMC

